\documentclass[10pt,a4paper]{amsart}
\usepackage[margin=19mm]{geometry}
\usepackage{amsmath,amssymb,mathtools}
\usepackage[scr=rsfs,cal=euler]{mathalfa}
\usepackage{xfrac}
\usepackage[unicode,hidelinks,bookmarks=false]{hyperref}
\newcommand{\ie}{i.e.\ }
\newcommand{\cf}{cf.\ }
\newcommand{\resp}{respectively\ }
\newcommand{\Subsection}{\subsection}
\providecommand{\nobreakdash}{\nobreak\hbox{-}}
\setlength{\emergencystretch}{2em}
\multlinegap0pt
\usepackage{bm}
\usepackage{bbm}
\usepackage{dsfont}
\usepackage{xspace}
\usepackage{cancel}
\usepackage{mathtools}
\usepackage{amsthm}
\makeatletter
\@ifundefined{Theorem}{\newtheorem{Theorem}{Theorem}}{}
\@ifundefined{Remark}{\newtheorem{Remark}[Theorem]{Remark}}{}
\makeatother
\title[The Buckling and Clamped Plate Problems on Differential Form]{The Buckling and Clamped Plate Problems on Differential Forms}
\author{Fida El Chami, Nicolas Ginoux, Georges Habib, Ola Makhoul, Simon Raulot}
\date{Source: arXiv:2412.05612v2 (CC BY 4.0)}
\begin{document}
\maketitle

\noindent\emph{Source and licence.} The Buckling and Clamped Plate Problems on Differential Forms. Version arXiv:2412.05612v2 (CC BY 4.0), distributed under the Creative Commons Attribution 4.0 International licence, as recorded in the arXiv licence metadata for this exact version and shown on the abstract page https://arxiv.org/abs/2412.05612v2. Full rights notice: licenses/arxiv-2412.05612-CC-BY-4.0.txt.\par\smallskip

\noindent\textit{Editorial scope.} Counted unit: the Theorem whose source label is \texttt{BuckAbs} in the article (source0.tex lines 633--646, source label \texttt{BuckAbs}), delivered together with the article's own complete original proof of it (source0.tex lines 648--691, one \texttt{proof} environment of 44 lines). No mathematics is written, repaired or re-derived: statement and proof are the article's own text line for line. Required context retained uncounted: Zcharacterization (lines 179--191); eq:laplacianG1 (lines 199--212); eq:laplacian (lines 214--216); eq:partialintegDelta2omega (lines 218--221); SymmetryBuck (lines 329--331); eq:Neumannpblonpforms (lines 582--588); VC\_Absolute (lines 591--593). Every definition, display and label of those lines is retained unchanged. Omitted from this package and not redistributed: 1--178, 192--198, 213--213, 217--217, 222--328, 332--581, 589--590, 594--632, 647--647, 692--1204. Nothing inside a retained excerpt is omitted or altered: every retained body is delivered exactly as the article's lines are written, so no omission receipt of any kind accompanies this package. Citations: the retained excerpts carry no citation command at all, so no bibliography entry of the article is transcribed and no reference is redistributed. Reference adaptation: none was needed and none was made; every reference inside the delivered unit resolves inside it, so no unresolved reference or citation is produced. Printed slips kept literal: no printed word, symbol, hypothesis or number of the counted statement or of its proof was altered. Layout and class: the article's own class is \texttt{sigma} with packages including ifpdf; the wrapper is the lane's pinned \texttt{amsart} editorial wrapper at 10pt on A4, which restates the theorem-like environments the retained text needs and adds the article's own macro definitions that the retained text needs (none), with extra wrapper packages (bm, bbm, dsfont, xspace, cancel, mathtools). The delivered numbering is the unit's own. Assets: no retained span contains a figure environment, a graphics-inclusion command, a PDF-page inclusion, a table or a TikZ picture; no image file is copied into this package, the package declares no asset dependency and no image file is redistributed with it. Licence: version arXiv:2412.05612v2 of this article is distributed under the Creative Commons Attribution 4.0 International licence, as recorded in the arXiv licence metadata for this exact version and shown on the abstract page https://arxiv.org/abs/2412.05612v2. A full rights notice is delivered at licenses/arxiv-2412.05612-CC-BY-4.0.txt. Characters: every retained excerpt is pure ASCII, so the delivered engine, XeLaTeX with its default Latin Modern text font, reports no missing character.
\begin{equation}\label{Zcharacterization}
\begin{aligned}
&\omega_{|\partial M} = 0,\\
&\nabla_\nu\omega_{|\partial M} = 0,
\end{aligned}
\quad\Longleftrightarrow\quad
\begin{aligned}
& \iota^*\omega = 0,\\
& \nu\lrcorner\omega = 0,\\
& \iota^*\delta\omega = 0,\\
& \nu\lrcorner\,{\rm d}\omega = 0.
\end{aligned}
\end{equation}

\begin{gather}
\int_M\bigl\langle\Delta\omega,\omega'\bigr\rangle {\rm d}\mu = \int_M\bigl(\bigl\langle {\rm d}\omega,{\rm d}\omega'\bigr\rangle +\bigl\langle\delta\omega,\delta\omega'\bigr\rangle\bigr) {\rm d}\mu \nonumber\\
\hphantom{\int_M\bigl\langle\Delta\omega,\omega'\bigr\rangle {\rm d}\mu =}{}
+\int_{\partial M}\bigl(\bigl\langle \nu\lrcorner {\rm d}\omega,\iota^*\omega'\bigr\rangle-\bigl\langle\iota^*\delta\omega,\nu\lrcorner\omega'\bigr\rangle\bigr) {\rm d}\sigma \label{eq:laplacianG1}\\
\hphantom{\int_M\bigl\langle\Delta\omega,\omega'\bigr\rangle {\rm d}\mu}{}
= \int_M\bigl\langle\omega,\Delta\omega'\bigr\rangle {\rm d}\mu + \int_{\partial M}\bigl(\bigl\langle\nu\lrcorner
{\rm d}\omega,\iota^*\omega'\bigr\rangle-\bigl\langle\iota^*\omega,\nu\lrcorner
{\rm d}\omega'\bigr\rangle
\nonumber\\
\hphantom{\int_M\bigl\langle\Delta\omega,\omega'\bigr\rangle {\rm d}\mu =\int_M\bigl\langle\omega,\Delta\omega'\bigr\rangle {\rm d}\mu + \int_{\partial M}}{}
+\bigl\langle\nu\lrcorner\omega,
\iota^*\delta\omega'\bigr\rangle-\bigl\langle\iota^*\delta\omega,
\nu\lrcorner\omega'\bigr\rangle\bigr){\rm d}\sigma,\label{eq:laplacianG}
\end{gather}

\begin{equation}\label{eq:laplacian}
\int_M\bigl\langle\Delta\omega,\omega'\bigr\rangle {\rm d}\mu=\int_M\bigl(\bigl\langle {\rm d}\omega,{\rm d}\omega'\bigr\rangle+\bigl\langle\delta\omega,\delta\omega'\bigr\rangle\bigr){\rm d}\mu=\int_M\bigl\langle\omega,\Delta\omega'\bigr\rangle {\rm d}\mu.
\end{equation}

\begin{align}
\nonumber\int_M\bigl\langle\Delta^2\omega,\omega'\bigr\rangle {\rm d}\mu ={}& \int_M\bigl\langle\Delta \omega,\Delta\omega'\bigr\rangle {\rm d}\mu+\int_{\partial M}\bigl(\bigl\langle\nu\lrcorner {\rm d}\Delta\omega,\iota^*\omega'\bigr\rangle-\bigl\langle\iota^*\Delta\omega,\nu\lrcorner {\rm d}\omega'\bigr\rangle\bigr){\rm d}\sigma\\
 &{}{+}\,\int_{\partial M}\bigl(\bigl\langle\nu\lrcorner\Delta\omega,\iota^*\delta\omega'\bigr\rangle-\bigl\langle\iota^*\delta\Delta\omega,\nu\lrcorner\omega'\bigr\rangle\bigr){\rm d}\sigma\label{eq:partialintegDelta2omega}
\end{align}

\begin{Remark}\label{SymmetryBuck}
When $M$ is oriented, the Hodge $\star$ operator is an isometry commuting with the Laplacian and preserving the boundary conditions in the buckling problem so that $\Lambda_{i,p}=\Lambda_{i,n-p}$ for any $i\geq 1$ and $1\leq p\leq n$.
\end{Remark}

\begin{equation}\label{eq:Neumannpblonpforms}
\begin{aligned}
&\Delta \omega =\mu \omega\ \quad &&\textrm{on } M, \\
&\nu\lrcorner\omega =0\ \quad&&\textrm{on }\partial M, \\
&\nu\lrcorner\,{\rm d}\omega = 0\ \quad &&\textrm{on }\partial M.
\end{aligned}
\end{equation}

\begin{equation}\label{VC_Absolute}
\mu_{1,p}=\inf_{\omega\in\Omega^p(M)\setminus\{0\}}\Biggl\{\frac{\|{\rm d}\omega\|_{L^2(M)}^2+\|\delta\omega\|_{L^2(M)}^2}{\|\omega\|_{L^2(M)}^2},\ \nu\lrcorner\omega=0\textrm{ and }\omega\in H_A^p(M)^\perp\Biggr\}.
\end{equation}

\begin{Theorem}\label{BuckAbs}
On an $n$-dimensional oriented compact Riemannian manifold $(M^n,g)$ with smooth boundary, we have
\begin{equation}\label{BuckAbs1}
\max(\mu_{1,p},\mu_{1,n-p})\leq\Lambda_{1,p}
\end{equation}
for all $p=0,\dots,n$. Moreover, equality occurs for some $p$ if and only if there exist $k\in\{p,n-p\}$, $\omega\in\Omega^k(M)$ as well as $\omega_0\in H_A^k(M)\setminus\{0\}$ satisfying the overdetermined boundary value problem%
\begin{equation}\label{ODP}
\begin{aligned}
&\Delta \omega =\mu_{1,k} \omega \ \quad&&\textrm{on } M, \\
&\omega = \omega_0 \ \quad&& \textrm{on }\partial M,\\
&\iota^*(\delta\omega) =0,\ \nu\lrcorner {\rm d}\omega =0 \ \quad&& \textrm{on }\partial M.
\end{aligned}
\end{equation}
\end{Theorem}

\begin{proof} Let $\omega_1\in\Omega^p(M)$ be an eigenform of the buckling problem associated to the eigenvalue~$\Lambda_{1,p}$.
We denote by $\omega_0$ the $L^2(M)$-orthogonal projection of $\omega_1$ onto \smash{$H_A^p(M)$} and we set $\omega:=\omega_1-\omega_0\in \Omega^p(M)$.
Note that $\omega_0$ may vanish, which is the case if and only if $\nu\lrcorner\,{\rm d}\Delta\omega_1$ is $L^2(\partial M)$-orthogonal to $\iota^*H_A^p(M)$ by~\eqref{eq:partialintegDelta2omega}.
Then ${\rm d}\omega={\rm d}\omega_1$ and $\delta\omega=\delta\omega_1$ on $M$, and in particular $\Delta\omega=\Delta\omega_1$, as well as $\nu\lrcorner\,\omega=0$ and $\nu\lrcorner\,{\rm d}\omega=0$ along $\partial M$.
For this form $\omega$, using the first equality in~(\ref{eq:laplacian}) and the Cauchy--Schwarz inequality leads to
\begin{gather*}
\bigl(\|{\rm d}\omega\|_{L^2(M)}^2+\|\delta\omega\|_{L^2(M)}^2\bigr)^2 = (\Delta\omega,\omega )_{L^2(M)}^2 \leq \|\Delta\omega\|_{L^2(M)}^2 \|\omega\|_{L^2(M)}^2,
\end{gather*}
so that, together with the fact that $\omega\neq0$ (otherwise, $\Lambda_{1,p}=0$),
\[
\frac{\|{\rm d}\omega\|_{L^2(M)}^2+\|\delta\omega\|_{L^2(M)}^2}{\|\omega\|_{L^2(M)}^2}\leq\frac{\|\Delta\omega\|_{L^2(M)}^2}{\|{\rm d}\omega\|_{L^2(M)}^2+\|\delta\omega\|_{L^2(M)}^2}.
\]
Note that \smash{$\|{\rm d}\omega\|_{L^2(M)}^2+\|\delta\omega\|_{L^2(M)}^2=\|{\rm d}\omega_1\|_{L^2(M)}^2+\|\delta\omega_1\|_{L^2(M)}^2>0$}.
Moreover, since $\Delta\omega=\Delta\omega_1$ and $\omega_1$ is an eigenform associated with $\Lambda_{1,p}$, we deduce that
\[
\frac{\|{\rm d}\omega\|_{L^2(M)}^2+\|\delta\omega\|_{L^2(M)}^2}{\|\omega\|_{L^2(M)}^2}\leq\frac{\|\Delta\omega_1\|_{L^2(M)}^2}{\|{\rm d}\omega_1\|_{L^2(M)}^2+\|\delta\omega_1\|_{L^2(M)}^2}=\Lambda_{1,p}.
\]
But $\nu\lrcorner\omega=0$ and \smash{$(\omega,\omega')_{L^2(M)}=0$} for all $\omega'\in H_A^p(M)$ so that the $p$-form $\omega$ can be taken as a~test-form in~(\ref{VC_Absolute}) and therefore $\displaystyle\mu_{1,p}\leq \Lambda_{1,p}$.
The same argument ensures that $\mu_{1,n-p}\leq\Lambda_{1,n-p}$ and the result follows from $\Lambda_{1,n-p}=\Lambda_{1,p}$ as was noticed in Remark~\ref{SymmetryBuck}.

If equality occurs, then it can be assumed without loss of generality that $\mu_{1,p}=\Lambda_{1,p}$ up to replacing $p$ by $n-p$.
Then the $p$-form $\omega$ is a $\mu_{1,p}$-Laplace-eigenform with respect to the absolute boundary condition, so that it satisfies the problem~(\ref{eq:Neumannpblonpforms}).
On the other hand, it is clear that, since $\omega=\omega_1-\omega_0$
%\footnote{Would $\omega_1-t\omega_0$ be better-suited?}
with $\omega_1$ a buckling-eigenform and $\omega_0\in H^p_A(M)$, we have $\iota^*\omega=-\iota^*\omega_0$ and ${\iota^*(\delta\omega)=0}$, therefore $\omega$ solves~\eqref{ODP} where $-\omega_0$ takes the role of $\omega_0$.
Note that necessarily the $\omega_0$ from~\eqref{ODP} cannot vanish, otherwise $\omega=0$ by the unique continuation property for second-order elliptic systems, which would be a contradiction.
This gives the first part of the equality case.
Conversely, assume the existence of $\omega\in\Omega^p(M)$ satisfying~(\ref{ODP}) with $\omega_0\in H^p_A(M)\setminus\{0\}$.
Let $\omega_1:=\omega-\omega_0\in\Omega^p(M)$.
Notice that $\omega_1\neq0$ (otherwise $\mu_{1,p}=0$) and that $\omega$ is $L^2(M)$-orthogonal to $H_A^p(M)$ by integrating $\Delta\omega=\mu_{1,p}\omega$ against any $\omega_0'\in H_A^p(M)$ on $M$ and using~\eqref{eq:laplacianG1}.
Then it is not difficult to check using~(\ref{Zcharacterization}) that
\[
\omega_{1}{}_{|_{\partial M}}=0\qquad\text{and}\qquad\nabla_\nu\omega_{1}{}_{|_{\partial M}}=0,
\]
so that $\omega_1$ can be taken as a test-form in the variational characterization of $\Lambda_{1,p}$.
On the other hand, since $\omega_0$ is a harmonic form, we have ${\rm d}\omega_1={\rm d}\omega$ and $\delta\omega_1=\delta\omega$ and then
\[
\Lambda_{1,p}\leq\frac{\|\Delta\omega_1\|^2_{L^2(M)}}{\|{\rm d}\omega_1\|^2_{L^2(M)}+\|\delta\omega_1\|^2_{L^2(M)}}=\frac{\|\Delta\omega\|^2_{L^2(M)}}{\|{\rm d}\omega\|^2_{L^2(M)}+\|\delta\omega\|^2_{L^2(M)}}=\mu_{1,p}.
\]
This implies that $\mu_{1,p}=\Lambda_{1,p}$.
Note that~\eqref{ODP} is invariant under the Hodge star operator: if~${\omega\in\Omega^p(M)}$ solves~\eqref{ODP}, then so does $\star\omega\in\Omega^{n-p}(M)$ for the same $\mu=\mu_{1,p}$.
This shows that, if $\mu_{1,p}=\Lambda_{1,p}$, then actually $\mu_{1,p}$ must also be a Laplace-eigenvalue with absolute boundary conditions on $n-p$-forms, however it does not have to be the smallest one.
This concludes the~proof.
\end{proof}

\end{document}
